%==============================================================================
\part{Realisation}
%==============================================================================

\section{A prototype}
\label{sec:proto}

\subsection{What the prototype is, and what it is not}

The prototype is a compiler and executor for \HFQ{} written in Python. It parses
plan text into the intermediate representation of \cref{def:plan}, performs the
static check of \cref{thm:static}, solves the allocation of
\cref{thm:allocation}, executes the plan against source adapters, assigns
verdicts by \cref{def:verdicts}, and writes the result --- verdicts, payloads,
retention figures, allocation, and provenance --- as a JSON document.

It is not a benchmark. Every adapter in the prototype resolves against a local
fixture or a local engine. No request leaves the machine. This is a deliberate
constraint and not a temporary limitation of the implementation: the properties
under test are properties of the compiler and the verdict assignment, and a
live service can neither confirm nor refute them while introducing a dependency
on a third party's availability, policy, and current snapshot. What a live
service would test --- whether real cross-reference tables have the retention
figures one hopes --- is a question this paper does not answer and does not
claim to (\cref{sec:limits}).

\subsection{Pipeline}

\begin{enumerate}[label=(\arabic*),leftmargin=2.4em]
\item \textbf{Parse.} Plan text to a list of steps
$(x_i, \Src_i, \rho_i, \beta_i, b_i)$, together with the plan budget. The
grammar is small: a plan is a budget declaration followed by \texttt{let}
bindings and a terminal \texttt{emit}.
\item \textbf{Resolve.} Each source name is resolved against a registry of
adapters, each of which declares $\Capset(\Src)$ as an explicit set of feature
symbols. The declaration is data, written by the adapter author, and is the
locus of the honesty assumption (\cref{rem:honesty-assumption}).
\item \textbf{Check.} Compute $\Req(\rho_i)$ for each step by structural
recursion over the abstract request, and test
$\Req(\rho_i) \subseteq \Capset(\Src_i)$. On failure the executor halts before
issuing any request and emits a refusal document naming the missing features
and the step --- \cref{cor:refuse-before-contact} made operational.
\item \textbf{Allocate.} Solve \cref{def:alloc} by bisection on the shadow price
$p$, using $\gamma_i(e) = w_i \log(1 + e)$ as the default yield, whose
derivative $w_i/(1+e)$ inverts in closed form to
$e = w_i/p - 1$. Steps with step-function yields are charged first
(\cref{rem:concavity-fails}).
\item \textbf{Execute.} In sequence order, apply (R1)--(R6). A step reaching
(R4) calls its adapter's lowering, which emits the canonical concrete request of
\cref{cons:longhand}; the adapter evaluates it against its fixture and returns a
result set.
\item \textbf{Emit.} Serialise to JSON.
\end{enumerate}

\subsection{Adapters}

Four adapter kinds are implemented, chosen to exercise the four source kinds of
\cref{obs:kind} rather than to be complete.

\begin{itemize}[leftmargin=2.4em]
\item A \emph{graph-pattern adapter} over a local RDF store
\citep{RDFLib,Pyoxigraph}, declaring
$\{\textsf{pattern}, \textsf{path}, \textsf{bind}, \textsf{filter}, \textsf{agg}\}$.
Its lowering emits longhand triple patterns and supplies bound sets through the
value-binding construct, never by concatenation.
\item A \emph{lookup adapter} over a local key--value fixture, declaring
$\{\textsf{lookup}, \textsf{link}\}$ and \emph{not} \textsf{pattern}. It is the
prototype's stand-in for a flat-file REST interface, and its restricted
declaration is what makes \cref{thm:static} fire on plans that ask it for a
join.
\item An \emph{ontology adapter} over a local class hierarchy, declaring
$\{\textsf{path}, \textsf{lookup}\}$ --- transitive closure over subsumption but
no filtering and no aggregation.
\item A \emph{map adapter} implementing translation steps
(\cref{def:map}), which is the only adapter that reports retention and
amplification.
\end{itemize}

\begin{remark}[The declarations are the experiment]
\label{rem:decl-experiment}
Because the static check is decided entirely by the declarations, and the
declarations are chosen by us, the prototype cannot demonstrate that the check
is \emph{right} about any real service. What it demonstrates is that the check
is \emph{decidable and cheap}, that refusal precedes contact, and that the
verdict assigned to a rejected plan names the missing feature rather than
reporting an empty result. Those are the claims of \cref{thm:static} and
\cref{cor:refuse-before-contact}, and they are properties of the compiler.
\end{remark}

\subsection{The output document}

The JSON document produced by an execution has one entry per step, and the
schema is fixed by the theory rather than by convenience:

\begin{lstlisting}[caption={Shape of a step entry in the emitted JSON. The \texttt{blocker} field is absent exactly when the verdict is \vempty{} or \vans{}, per \cref{def:blocker}.},label=lst:json,language=,keywordstyle=]
{
  "step": "kids",
  "source": "CHEBI->KEGG",
  "verdict": "starved",
  "blocker": "corpus",
  "diagnosis": {
    "named_predecessor": "ecs",
    "reason": "retention 0.31 below declared expectation 0.6"
  },
  "retention": 0.31,
  "amplification": 1.14,
  "budget": {"allocated": 60, "spent": 60, "shadow_price": 0.0142},
  "payload": null,
  "provenance": {"snapshot": "fixture-v3", "lowered_form": "canonical"}
}
\end{lstlisting}

Three features of \cref{lst:json} are consequences of results above rather than
design taste. The \texttt{blocker} field is partial, because
\cref{def:blocker} is partial and $\vempty$ has no blocker. The
\texttt{payload} is \texttt{null} for every non-$\vans$ verdict, because
\cref{prin:verdict} says a payload accompanies only $\vans$ --- and an
implementation that returned a partial payload alongside a failure verdict
would reintroduce exactly the ambiguity \cref{cor:onebit} identifies.
\texttt{retention} and \texttt{amplification} are recorded separately, because
by \cref{prop:cardinality-uninformative} their product --- which is all the
output cardinality reveals --- determines neither.

\section{Validation design}
\label{sec:validation}

We enumerate the checks the prototype performs, in the interest of making it
clear which claims of this paper are exercised and which are not. Each is
executed against local fixtures.

\begin{enumerate}[label=(V\arabic*),leftmargin=3em]
\item \textbf{Static check decides without contact.} For a family of plans,
half well-capability and half not, verify that the ill-capability plans halt
with a refusal document, that the refusal names the correct missing features,
and that the adapter request counter is zero.
(\cref{thm:static}, \cref{cor:refuse-before-contact}.)
\item \textbf{Check cost is linear.} Measure the check's operation count
against $m|\Feat|$ over plans of growing length. (\cref{thm:static}(a).)
\item \textbf{Six verdicts are realised.} Construct the six configurations of
the proof of \cref{thm:six} --- each differing from a common baseline in one
respect --- and verify that the executor assigns six distinct verdicts.
(\cref{thm:six}.)
\item \textbf{The one-bit interface conflates four.} Run the same six
configurations through a deliberately impoverished adapter that returns only a
result set, and verify that four of the six are indistinguishable in its
output. (\cref{cor:onebit}.)
\item \textbf{Starvation is unreachable when $m=1$.} Over all single-step
configurations in the fixture family, verify that $\vstarve$ is never assigned
and $\bcorp$ never emitted; then verify that the predecessor perturbation of
\cref{prop:starve-reachable}(b) produces it in a two-step plan.
\item \textbf{Blame terminates.} For every starved execution, follow the
diagnosis chain and verify it terminates at a non-$\vstarve$ step within $m$
hops. (\cref{prop:blame}.)
\item \textbf{Retention factorises.} Over random partial maps, verify the
multiplicative identity of \cref{thm:retention}(a) exactly, and verify that the
measured surviving fraction lies within the bounds of (b) and (c). Record how
often each bound is attained and how wide the gap is
(\cref{rem:bounds-gap}).
\item \textbf{Cardinality is uninformative.} Realise the two maps of
\cref{prop:cardinality-uninformative} and verify equal output cardinality with
retentions differing by the predicted factor.
\item \textbf{Pairwise coverage does not determine chain coverage.} Realise the
two families of \cref{prop:pairwise-insufficient} and verify identical pairwise
retentions with differing end-to-end surviving fractions.
\item \textbf{Reordering does not recover coverage.} For plans admitting both
orders, verify equal surviving fractions and unequal request counts.
(\cref{thm:reorder}, \cref{rem:reorder-cost}.)
\item \textbf{Routes diverge without contradiction.} Construct parallel sound
routes with partial maps, verify that neither contains an element contradicting
the other, that the symmetric difference is non-empty, and that making every
map total collapses the divergence to zero.
(\cref{thm:route-extent}.)
\item \textbf{Allocation satisfies KKT.} Solve \cref{def:alloc} numerically and
verify that the marginal yields of all supported steps agree to tolerance, that
unsupported steps satisfy $\gamma_i'(0) \leq p^\ast$, and that the budget binds.
(\cref{thm:allocation}.)
\item \textbf{Sufficient budget is not sufficient.} Realise the two-step
counterexample of \cref{prop:necessary-not-sufficient} and verify a $\vref$
verdict despite $\Bud \geq \sum c_i$.
\item \textbf{Lowering is canonical.} For every abstract request in the fixture
family, verify that the emitted concrete form contains no predicate-object list
and no interpolated identifier outside a value-binding construct, and that two
abstract requests with the same denotation lower to the same string.
(\cref{cons:longhand}, \cref{thm:interpolation}.)
\item \textbf{Re-execution is stable.} Execute each plan twice over an unchanged
fixture with sufficient budget and verify identical verdicts and payloads; then
re-execute with a budget below one step's realised cost and verify that the
difference is confined to that step and its successors.
(\cref{thm:idempotent}.)
\end{enumerate}

We record what these do \emph{not} establish, since the enumeration invites the
opposite reading. (V1)--(V15) test the compiler and the executor against the
theory. They do not test the theory against the world. No check above would
detect a dishonest capability declaration, an inaccurate cross-reference table,
or a deployment that misinterprets the canonical form --- and the first and
third of those are precisely the failure modes \cref{rem:honesty-assumption}
and \cref{rem:interp-scope} identify as beyond the reach of the construction.

\section{Limitations}
\label{sec:limits}

\paragraph{Honesty is assumed, never verified.} The whole of \cref{thm:static}
and the greater part of \cref{sec:verdicts} rest on the assumption that a
source's declared capability set is honest (\cref{def:honest}). Under-declaration
costs expressiveness and is safe; over-declaration is unsound and, worse,
invisible --- an over-declared source accepts a request it will silently degrade,
and the degraded answer arrives as $\vans$ with a payload. We do not know how to
verify a declaration without probing the service, and probing a third party's
public service to characterise its internals is not something we undertake.
This is the single largest gap in the development and we do not minimise it.

\paragraph{Yields are posited, not estimated.} \Cref{thm:allocation} assumes
strictly concave differentiable yield functions and the prototype supplies
$w_i\log(1+e)$ by default. We have not estimated a yield function for any real
source, and the weights $w_i$ are configuration. An allocation optimal for the
wrong yields is not optimal. Furthermore \cref{rem:concavity-fails} excludes
all-or-nothing steps from the optimisation entirely, so the shadow price
governs a subset of the plan.

\paragraph{No regret bound for replanning.} \Cref{rem:replanning} adopts
re-solving after each step, and we have not bounded its total yield against the
clairvoyant allocation that knew the realised cardinalities in advance. This is
the natural theoretical question raised by
\cref{prop:necessary-not-sufficient} and we do not answer it.

\paragraph{Retention laws are unmeasured.} \Cref{thm:retention} is arithmetic
and holds of any partial maps. Whether real cross-reference tables have
retentions making \cref{prop:starve-common} bite, and whether the losses of
successive maps are correlated or independent, is an empirical question we
raise and do not settle. We note only that the independence hypothesis is
false in the direction that makes the bound conservative.

\paragraph{Namespace disjointness is a modelling decision.} \Cref{rem:disjoint}
records that we take namespaces to be pairwise disjoint. Real identifier spaces
overlap and are re-used, and a shared identifier appearing in two namespaces
would be conflated by the model. We do not know how badly.

\paragraph{Wall-clock time is not modelled.} All costs count requests
(\cref{sec:intro}, conventions). A plan whose requests are few and slow and one
whose requests are many and fast are indistinguishable to the allocator. The
verdict $\vtime$ mentions a clock but the budget does not, and reconciling the
two would require a second resource dimension the model does not carry.

\paragraph{The static check is not a completeness result.} \Cref{thm:static}
says a well-capability plan compiles. It does not say an ill-capability plan is
unanswerable: a source might answer a request outside its declared capabilities
through a route the adapter does not expose. Refusal is conservative, and
conservatism is the safe direction, but it is not optimality.

\paragraph{Structural interpolation is a smaller claim than correctness.}
\Cref{rem:interp-scope} states this in place and we repeat it here because the
result is the one most liable to be over-read. \Cref{thm:interpolation}
eliminates a family of defects by removing the author's access to the request
string. It does not make the request correct, and a deployment that misreads the
canonical form is untouched by it.

\paragraph{Nothing has been run against a live service.} By construction. The
consequence is that every quantitative figure in the prototype's output is a
figure about a fixture, and no claim in this paper is supported by measurement
of a public biological database.

\section{Conclusion}
\label{sec:conclusion}

The script of \cref{lst:script} is not bad code. It is the best code available
under a division of labour in which control flow lives in a host language and
data access lives in a query language, with no channel between them. Its five
defects are consequences of that division and not of its author.

We have described a language that moves the division. Control flow and data
access live together in a plan; the plan author never writes a request; requests
are generated by adapters from abstract steps and source declarations. The
consequences follow from that single relocation. Because requests are generated,
their spelling is canonical and a family of interpolation defects cannot occur
(\cref{thm:interpolation}). Because sources declare capabilities, compilability
is decidable before contact and refusal precedes any request
(\cref{thm:static}). Because a step returns a verdict rather than rows, the four
outcomes a row-returning interface conflates become distinguishable
(\cref{cor:onebit}) --- including one, starvation, that has no single-database
analogue and becomes reachable only in federation
(\cref{prop:starve-reachable}). Because translation is a first-class step, its
retention is measured rather than assumed, and the chain laws say what a plan
can and cannot infer about its own coverage
(\cref{thm:retention}, \cref{prop:pairwise-insufficient}) --- including that
reordering will not help (\cref{thm:reorder}) and that two sound routes may
disagree without either being wrong (\cref{thm:route-extent}). Because the
budget is declared, allocation is an optimisation with a single shadow price
rather than a priority list (\cref{thm:allocation}), though a sufficient total
is still not sufficient (\cref{prop:necessary-not-sufficient}).

What the development does not have is a way to check the one thing it most
depends on. A capability declaration is an assertion by an adapter author, and
an over-declaration is unsound and silent. Every static guarantee in this paper
is conditional on that assertion. Whether the condition can be discharged
without probing the services in question --- which is not ours to do --- is the
question we would most like answered and the one we leave open.
